\documentclass[reqno]{amsart}
\usepackage[T1]{fontenc}
\usepackage{amsmath,amssymb,amsthm}
\usepackage{xr-hyper}
\usepackage[hidelinks]{hyperref}
\externaldocument[so-]{20261124T220500--sign-only-proof-for-curves__workbook}

% Chapter number; \ref{chap} gives it.
\newcounter{chap}
% Chapter numbers, generated from 20261125T012000--grh-workbook__guide.org
% by scripts/make-chapters.py.  Do not edit: change the order of the
% index and rerun the script (./preview-site.sh and CI do this).
%
% Usage in a sheet, after \newcounter{chap}:
%   % Chapter numbers, generated from 20261125T012000--grh-workbook__guide.org
% by scripts/make-chapters.py.  Do not edit: change the order of the
% index and rerun the script (./preview-site.sh and CI do this).
%
% Usage in a sheet, after \newcounter{chap}:
%   % Chapter numbers, generated from 20261125T012000--grh-workbook__guide.org
% by scripts/make-chapters.py.  Do not edit: change the order of the
% index and rerun the script (./preview-site.sh and CI do this).
%
% Usage in a sheet, after \newcounter{chap}:
%   \input{chapters}\chapnumber{<sheet ID>}
\makeatletter
\providecommand\chapters@set[2]{\@namedef{chapters@#1}{#2}}
\chapters@set{20261125T001000}{1}% zeros-to-ramp-sums
\chapters@set{20261125T001100}{2}% inert-primes-and-estimate-b
\chapters@set{20261125T005000}{3}% galois-fourier-and-averaging
\chapters@set{20261125T005500}{4}% columns-from-a-sphere
\chapters@set{20261125T000000}{5}% slope-method
\chapters@set{20261126T030100}{6}% a-twist-for-each-prime
\chapters@set{20261125T003000}{7}% the-chain-and-the-square-root
\chapters@set{20261125T010000}{8}% rank-lemma-and-zero-estimate
\chapters@set{20261125T010500}{9}% poincare-series-and-near-orthogonality
\chapters@set{20261124T223000}{10}% elliptic-curves-for-the-proof
\chapters@set{20261125T009500}{11}% fock-model-and-theta-intertwiners
\chapters@set{20261124T223100}{12}% complex-multiplication-and-inert-reduction
\chapters@set{20261125T008000}{13}% one-cm-determinant
\chapters@set{20261124T214500}{14}% supersingular-theta-divisibility
\chapters@set{20261124T214600}{15}% the-isotropic-plane
\chapters@set{20261125T013700}{16}% pairings-on-a-formal-group-of-height-two
\chapters@set{20261125T013500}{17}% hecke-nondegeneracy-on-the-tree
\chapters@set{20261126T030200}{18}% the-bound-is-attained
\chapters@set{20261125T009000}{19}% kronecker-limit-and-chowla-selberg
\chapters@set{20261125T002000}{20}% definite-theta-lifts
\chapters@set{20261125T015500}{21}% siegel-weil-in-the-smallest-case
\chapters@set{20261125T002100}{22}% definite-theta-lifts-and-norms
\chapters@set{20261125T016500}{23}% hecke-operators-and-theta-lifts
\chapters@set{20261125T016700}{24}% divisibility-along-the-chain
\chapters@set{20261125T004000}{25}% denominators-and-sha
\chapters@set{20261125T004100}{26}% tate-twists-and-sha
\chapters@set{20261126T020430}{27}% the-minus-part-by-gauss-sums
\chapters@set{20261125T014100}{28}% congruence-modules
\chapters@set{20261125T014500}{29}% eisenstein-congruences-and-selmer-classes
\chapters@set{20261125T014700}{30}% pseudocharacters-and-the-length-bound
\chapters@set{20261125T011000}{31}% galois-representations-black-box
\chapters@set{20261126T030300}{32}% the-data
\chapters@set{20261126T020210}{33}% bessel-poincare-series-and-their-gram-operators
\chapters@set{20261126T030400}{34}% the-seesaw-identity
\chapters@set{20261126T030500}{35}% divisibility-of-the-composite-lift
\chapters@set{20261126T030600}{36}% theorem-a-short-forms
\chapters@set{20261126T030700}{37}% the-determinant-of-a-chain
\chapters@set{20261126T030800}{38}% the-cost-of-one-lift
\chapters@set{20261126T030900}{39}% denominators-of-projections-to-lifts
\chapters@set{20261126T031000}{40}% twists-with-few-congruences
\chapters@set{20261126T031100}{41}% theorem-c-and-the-main-inequality
\providecommand\chapnumber[1]{%
  \@ifundefined{chapters@#1}%
    {\PackageError{chapters}{Sheet #1 is not listed in the index}%
      {List it in 20261125T012000--grh-workbook__guide.org\MessageBreak
       and run scripts/make-chapters.py.}}%
    {\edef\thechap{\@nameuse{chapters@#1}}}}
\makeatother
\chapnumber{<sheet ID>}
\makeatletter
\providecommand\chapters@set[2]{\@namedef{chapters@#1}{#2}}
\chapters@set{20261125T001000}{1}% zeros-to-ramp-sums
\chapters@set{20261125T001100}{2}% inert-primes-and-estimate-b
\chapters@set{20261125T005000}{3}% galois-fourier-and-averaging
\chapters@set{20261125T005500}{4}% columns-from-a-sphere
\chapters@set{20261125T000000}{5}% slope-method
\chapters@set{20261126T030100}{6}% a-twist-for-each-prime
\chapters@set{20261125T003000}{7}% the-chain-and-the-square-root
\chapters@set{20261125T010000}{8}% rank-lemma-and-zero-estimate
\chapters@set{20261125T010500}{9}% poincare-series-and-near-orthogonality
\chapters@set{20261124T223000}{10}% elliptic-curves-for-the-proof
\chapters@set{20261125T009500}{11}% fock-model-and-theta-intertwiners
\chapters@set{20261124T223100}{12}% complex-multiplication-and-inert-reduction
\chapters@set{20261125T008000}{13}% one-cm-determinant
\chapters@set{20261124T214500}{14}% supersingular-theta-divisibility
\chapters@set{20261124T214600}{15}% the-isotropic-plane
\chapters@set{20261125T013700}{16}% pairings-on-a-formal-group-of-height-two
\chapters@set{20261125T013500}{17}% hecke-nondegeneracy-on-the-tree
\chapters@set{20261126T030200}{18}% the-bound-is-attained
\chapters@set{20261125T009000}{19}% kronecker-limit-and-chowla-selberg
\chapters@set{20261125T002000}{20}% definite-theta-lifts
\chapters@set{20261125T015500}{21}% siegel-weil-in-the-smallest-case
\chapters@set{20261125T002100}{22}% definite-theta-lifts-and-norms
\chapters@set{20261125T016500}{23}% hecke-operators-and-theta-lifts
\chapters@set{20261125T016700}{24}% divisibility-along-the-chain
\chapters@set{20261125T004000}{25}% denominators-and-sha
\chapters@set{20261125T004100}{26}% tate-twists-and-sha
\chapters@set{20261126T020430}{27}% the-minus-part-by-gauss-sums
\chapters@set{20261125T014100}{28}% congruence-modules
\chapters@set{20261125T014500}{29}% eisenstein-congruences-and-selmer-classes
\chapters@set{20261125T014700}{30}% pseudocharacters-and-the-length-bound
\chapters@set{20261125T011000}{31}% galois-representations-black-box
\chapters@set{20261126T030300}{32}% the-data
\chapters@set{20261126T020210}{33}% bessel-poincare-series-and-their-gram-operators
\chapters@set{20261126T030400}{34}% the-seesaw-identity
\chapters@set{20261126T030500}{35}% divisibility-of-the-composite-lift
\chapters@set{20261126T030600}{36}% theorem-a-short-forms
\chapters@set{20261126T030700}{37}% the-determinant-of-a-chain
\chapters@set{20261126T030800}{38}% the-cost-of-one-lift
\chapters@set{20261126T030900}{39}% denominators-of-projections-to-lifts
\chapters@set{20261126T031000}{40}% twists-with-few-congruences
\chapters@set{20261126T031100}{41}% theorem-c-and-the-main-inequality
\providecommand\chapnumber[1]{%
  \@ifundefined{chapters@#1}%
    {\PackageError{chapters}{Sheet #1 is not listed in the index}%
      {List it in 20261125T012000--grh-workbook__guide.org\MessageBreak
       and run scripts/make-chapters.py.}}%
    {\edef\thechap{\@nameuse{chapters@#1}}}}
\makeatother
\chapnumber{<sheet ID>}
\makeatletter
\providecommand\chapters@set[2]{\@namedef{chapters@#1}{#2}}
\chapters@set{20261125T001000}{1}% zeros-to-ramp-sums
\chapters@set{20261125T001100}{2}% inert-primes-and-estimate-b
\chapters@set{20261125T005000}{3}% galois-fourier-and-averaging
\chapters@set{20261125T005500}{4}% columns-from-a-sphere
\chapters@set{20261125T000000}{5}% slope-method
\chapters@set{20261126T030100}{6}% a-twist-for-each-prime
\chapters@set{20261125T003000}{7}% the-chain-and-the-square-root
\chapters@set{20261125T010000}{8}% rank-lemma-and-zero-estimate
\chapters@set{20261125T010500}{9}% poincare-series-and-near-orthogonality
\chapters@set{20261124T223000}{10}% elliptic-curves-for-the-proof
\chapters@set{20261125T009500}{11}% fock-model-and-theta-intertwiners
\chapters@set{20261124T223100}{12}% complex-multiplication-and-inert-reduction
\chapters@set{20261125T008000}{13}% one-cm-determinant
\chapters@set{20261124T214500}{14}% supersingular-theta-divisibility
\chapters@set{20261124T214600}{15}% the-isotropic-plane
\chapters@set{20261125T013700}{16}% pairings-on-a-formal-group-of-height-two
\chapters@set{20261125T013500}{17}% hecke-nondegeneracy-on-the-tree
\chapters@set{20261126T030200}{18}% the-bound-is-attained
\chapters@set{20261125T009000}{19}% kronecker-limit-and-chowla-selberg
\chapters@set{20261125T002000}{20}% definite-theta-lifts
\chapters@set{20261125T015500}{21}% siegel-weil-in-the-smallest-case
\chapters@set{20261125T002100}{22}% definite-theta-lifts-and-norms
\chapters@set{20261125T016500}{23}% hecke-operators-and-theta-lifts
\chapters@set{20261125T016700}{24}% divisibility-along-the-chain
\chapters@set{20261125T004000}{25}% denominators-and-sha
\chapters@set{20261125T004100}{26}% tate-twists-and-sha
\chapters@set{20261126T020430}{27}% the-minus-part-by-gauss-sums
\chapters@set{20261125T014100}{28}% congruence-modules
\chapters@set{20261125T014500}{29}% eisenstein-congruences-and-selmer-classes
\chapters@set{20261125T014700}{30}% pseudocharacters-and-the-length-bound
\chapters@set{20261125T011000}{31}% galois-representations-black-box
\chapters@set{20261126T030300}{32}% the-data
\chapters@set{20261126T020210}{33}% bessel-poincare-series-and-their-gram-operators
\chapters@set{20261126T030400}{34}% the-seesaw-identity
\chapters@set{20261126T030500}{35}% divisibility-of-the-composite-lift
\chapters@set{20261126T030600}{36}% theorem-a-short-forms
\chapters@set{20261126T030700}{37}% the-determinant-of-a-chain
\chapters@set{20261126T030800}{38}% the-cost-of-one-lift
\chapters@set{20261126T030900}{39}% denominators-of-projections-to-lifts
\chapters@set{20261126T031000}{40}% twists-with-few-congruences
\chapters@set{20261126T031100}{41}% theorem-c-and-the-main-inequality
\providecommand\chapnumber[1]{%
  \@ifundefined{chapters@#1}%
    {\PackageError{chapters}{Sheet #1 is not listed in the index}%
      {List it in 20261125T012000--grh-workbook__guide.org\MessageBreak
       and run scripts/make-chapters.py.}}%
    {\edef\thechap{\@nameuse{chapters@#1}}}}
\makeatother
\chapnumber{20261124T221500}
\renewcommand{\thesection}{\thechap.\arabic{section}}

\theoremstyle{definition}
\newtheorem{problem}{Problem}[section]
\theoremstyle{plain}
\newtheorem{theorem}{Theorem}[section]

\newcommand{\Q}{\mathbb{Q}}
\newcommand{\C}{\mathbb{C}}
\newcommand{\F}{\mathbb{F}}
\newcommand{\A}{\mathbb{A}}
\newcommand{\G}{\mathbb{G}}
\newcommand{\Ql}{\overline{\Q}_\ell}
\newcommand{\cF}{\mathcal{F}}
\newcommand{\Fr}{\mathrm{Fr}}
\newcommand{\Tr}{\operatorname{Tr}}

\setlength{\emergencystretch}{2em}

\title{Deligne's bound by the sign-only method}
\author{}
\date{November 24, 2026}

\begin{document}
\maketitle
\refstepcounter{chap}\label{chap}

\paragraph{Prerequisites.}
Chapter~\ref{so-chap}, \emph{The sign-only proof of the Riemann hypothesis for curves}, whose notation, background and results are used throughout.  Constructible \(\ell\)-adic sheaves \cite[Rapport]{SGA45}, \cite{Mil}.  Noether normalization, and generic smoothness of a reduced scheme of finite type over a perfect field.

\section{The lowest degree}\label{sec:lowest}

Nothing in \S\ref{so-sec:curves} used \(\dim U = 1\) except the finite map \(\pi\) to \(\A^1\) and the shape of the coupling \(\pi(u_i)\pi(u_{i+1})\).  What does the sign-only method give for \(U\) of dimension \(d\)?  Fix \(\iota \colon \Ql \to \C\).  A constructible sheaf \(\cF\) on \(X\) has \emph{weights at most} \(w\) if at every closed point \(x\) all eigenvalues of the local Frobenius have absolute value at most \(\lvert \kappa(x) \rvert^{w/2}\).  It is \emph{pure of weight} \(w\) if all these absolute values equal \(\lvert \kappa(x) \rvert^{w/2}\).

\begin{theorem}[Deligne]\label{th:deligne}
Let \(X\) be separated of finite type over \(\F_q\), and \(\cF\) a constructible \(\Ql\)-sheaf on \(X\).
\begin{itemize}
\item[(i)] If \(\cF\) has weights at most \(w\), every eigenvalue of \(\Fr\) on \(H^i_c(X, \cF)\) has absolute value at most \(q^{(w+i)/2}\).
\item[(ii)] If \(X\) is smooth and proper and \(\cF\) is lisse and pure of weight \(w\), every eigenvalue of \(\Fr\) on \(H^i(X, \cF)\) has absolute value \(q^{(w+i)/2}\).
\end{itemize}
\end{theorem}

These are theorems of Deligne \cite{Del74}, \cite[Th\'eor\`eme~3.3.1]{Del80}.  By a result of Laumon and Deligne's dévissage, (i) follows from the case of a lisse sheaf on an open subset of \(\A^1\) \cite[Remark~1.3]{Zh}, and that case is written out by the sign-only method in \cite{Zh}.  The proof below instead treats every dimension directly.  It is that of \cite{Ald}, and as recorded in Section~\ref{so-sec:identity}, no human has checked that text.  Zhang states that the argument adapts to this setting, but he does not write it out \cite[Remark~1.3]{Zh}.

\begin{problem}\label{pr:affine}
For \(X = \A^d\) and \(\cF = \Ql\), \(H^i_c(X, \Ql) = 0\) for \(i \ne 2d\), and \(\Fr\) acts on the line \(H^{2d}_c(X, \Ql)\) by \(q^d\).  So (i) is attained for \(w = 0\), \(i = 2d\).
\end{problem}

Let \(U\) be smooth and affine of pure dimension \(d \ge 1\) over \(\F_q\), and \(\cF\) lisse on \(U\) with \(\lvert \iota\Tr(\Fr_u \mid \cF_{\bar u}) \rvert \le C q^{nw/2}\) for all \(n\) and \(u \in U(\F_{q^n})\).  Let \(\pi \colon U \to \A^d\) be finite (Noether normalization), \(B\) the number of generators of \(\mathcal O(U)\) as an \(\mathcal O(\A^d)\)-module, and \(z_1, \dots, z_s\) generators of the \(\F_q\)-algebra \(\mathcal O(U)\).  Let \(b\) be a nondegenerate bilinear form on \(\A^d\), and let \(Y_{a,c} \subset U^k\) be the zero set of
\[
  \sum_{i=1}^{k-1} b\bigl(\pi(u_i), \pi(u_{i+1})\bigr) + \sum_{i=1}^{k} \sum_{j=1}^{s} a_{ij}\, z_j(u_i) + c .
\]
This is the note's ``vectorize to go up dimensions'' \cite[footnote~3]{Alp}.  Let \(X := U^k\), \(D := kd\), \(\cF_k := \cF^{\boxtimes k}\), and define \(M_n\), \(N_n\), \(R_n\) and \(K\) for \(Y = Y_{a,c}\) as in \S\ref{so-sec:curves}, with \(H^{D-1}_c\) and \(H^{D+1}_c\) in the definition of \(K\).

\begin{problem}\label{pr:vector-d}
Let \(E\) be a finite field with \(r\) elements, \(\psi\) a nontrivial additive character of \(E\), \(b_1, \dots, b_{k-1}\) nondegenerate bilinear forms on \(E^d\), and \(t_1, \dots, t_k \colon E^d \to \C\) bounded by \(T\).  For \(h \in E^\times\) and \(c \in E\),
\[
  \Bigl\lvert \sum_{x \in (E^d)^k} t_1(x_1) \cdots t_k(x_k)\, \psi\Bigl(h \sum_{i=1}^{k-1} b_i(x_i, x_{i+1}) + hc\Bigr) \Bigr\rvert \le T^k r^{d(k+1)/2}.
\]
\end{problem}

This is \cite[Lemma~3]{Ald}.

\begin{problem}\label{pr:Gamma-d}
The map \(u \mapsto \bigl(z(u_1), \dots, z(u_k), \sum_i b(\pi(u_i), \pi(u_{i+1}))\bigr)\) is a closed immersion \(U^k \to \A^{ks+1}\), and the \(Y_{a,c}\) are its hyperplane sections with last coefficient \(1\).
\end{problem}

\begin{problem}\label{pr:fibres}
Over every finite extension of \(\F_q\), each fibre of \(\pi\) has at most \(B\) rational points.
\end{problem}

\begin{problem}\label{pr:kunneth-d}
\(H^j_c(X, \cF_k) = 0\) for \(j < D\), and \(H^D_c(X, \cF_k) = H^d_c(U, \cF)^{\otimes k}\).
\end{problem}

By Problems~\ref{so-pr:wl-affine} and~\ref{pr:Gamma-d} with \(\mathcal{G} = \cF_k^\vee\), there are \(e \ge 1\) and \((a, c)\) over \(\F_r\), \(r := q^e\), for which \(Y := Y_{a,c}\) has property (WL) for \(\cF_k\).

\begin{problem}\label{pr:prop4-d}
Let \(\Fr\) be the geometric Frobenius over \(\F_r\).  For all \(n \ge 1\),
\(
  (-1)^{D-1} R_n = \Tr\bigl(\Fr^n \mid H^d_c(U, \cF)^{\otimes k}\bigr) + \Tr\bigl(\Fr^n \mid K\bigr).
\)
\end{problem}

\begin{problem}\label{pr:R-d}
\(\lvert R_n \rvert \le (BC)^k\, r^{n(k(w+d)/2 + d/2 + 1)}\).
\end{problem}

Problems~\ref{pr:prop4-d}, \ref{pr:R-d} and~\ref{so-pr:lemma2}, the latter with its \(a\) equal to \((w+d)/2\) and \(c_k = d/2 + 1\), prove:

\begin{theorem}[{\cite[Theorem~9]{Ald}}]\label{th:lowest}
Every eigenvalue of \(\Fr\) on \(H^d_c(U, \cF)\) has absolute value at most \(q^{(w+d)/2}\).
\end{theorem}

Below degree \(d\) there is nothing to prove, since \(H^i_c(U, \cF) = 0\) for \(i < d\).  The degrees above \(d\) are reached by induction on the dimension.

\section{Induction on the dimension}\label{sec:induction}

\begin{problem}\label{pr:gm}
\(H^1_c(\G_m, \Ql)\) is a line on which \(\Fr\) acts by \(1\), and \(H^2_c(\G_m, \Ql)\) a line on which it acts by \(q\).  So the inequality in (i) can be strict.
\end{problem}

\begin{problem}\label{pr:dim0}
Theorem~\ref{th:deligne}(i) holds when \(X\) has dimension \(0\).
\end{problem}

\begin{problem}\label{pr:extension}
A bound \(\lvert \alpha \rvert \le q^{(w+i)/2}\) for the eigenvalues of \(\Fr\) on \(H^i_c(X, \cF)\) may be proved after replacing \(\F_q\) by \(\F_{q^e}\).
\end{problem}

\begin{problem}\label{pr:devissage}
Let \(U \subset X\) be open with closed complement \(Z\).  Every eigenvalue of \(\Fr\) on \(H^i_c(X, \cF)\) is an eigenvalue on \(H^i_c(U, \cF)\) or on \(H^i_c(Z, \cF)\).
\end{problem}

\begin{problem}\label{pr:lisse-open}
Every \(X\) of dimension \(d\) has a dense open subscheme \(U\) such that \(U_{\mathrm{red}}\) is a finite disjoint union of smooth affine irreducible schemes and \(\cF\) is lisse on \(U\).
\end{problem}

\begin{problem}\label{pr:traces}
If \(\cF\) is lisse of rank \(\rho\) on \(U\) with weights at most \(w\), then \(\lvert \iota\Tr(\Fr_u \mid \cF_{\bar u}) \rvert \le \rho\, q^{nw/2}\) for all \(n\) and \(u \in U(\F_{q^n})\).
\end{problem}

\begin{problem}\label{pr:higher}
Let \(X\) be smooth and affine of pure dimension \(d \ge 1\), \(\cF\) lisse, and \(i > d\).  Let \(Y \subset X\) be a hyperplane section with property (WL) for \(\cF\), defined over \(\F_q\) (after replacing \(\F_q\) by a finite extension).  If the eigenvalues of \(\Fr\) on \(H^{i-2}_c(Y, \cF)\) have absolute value at most \(q^{(w+i-2)/2}\), then those on \(H^i_c(X, \cF)\) have absolute value at most \(q^{(w+i)/2}\).
\end{problem}

Problems~\ref{pr:dim0}--\ref{pr:higher} and Theorem~\ref{th:lowest}, with Problem~\ref{so-pr:wl-affine} applied to an embedding of \(X\) to produce the hyperplane sections of Problem~\ref{pr:higher}, prove Theorem~\ref{th:deligne}(i) by induction on \(\dim X\) \cite[\S 8]{Ald}.

\begin{problem}\label{pr:purity}
Theorem~\ref{th:deligne}(i) implies Theorem~\ref{th:deligne}(ii).
\end{problem}

\section*{Hints}

\subsection*{\S\ref{sec:lowest}}
\begin{description}
\item[\ref{pr:affine}] The K\"unneth formula, and \(H^*_c(\A^1, \Ql) = \Ql(-1)\) in degree \(2\).
\item[\ref{pr:vector-d}] As in Problem~\ref{so-pr:parseval-k}; the matrix \(\bigl(\psi(h\, b(x,y))\bigr)_{x,y \in E^d}\) times its adjoint is \(r^d\), as in Problem~\ref{so-pr:unitary}.
\item[\ref{pr:Gamma-d}] As in Problem~\ref{so-pr:Gamma}.
\item[\ref{pr:fibres}] A fibre is the spectrum of an algebra spanned by \(B\) elements.
\item[\ref{pr:kunneth-d}] The K\"unneth formula and the vanishing of \(H^j_c(U, \cF)\) for \(j < d\).
\item[\ref{pr:prop4-d}] As in Problem~\ref{so-pr:prop4}, with Problems~\ref{so-pr:gysin} and~\ref{pr:kunneth-d}.
\item[\ref{pr:R-d}] As in Problems~\ref{so-pr:char-sum} and~\ref{so-pr:R-curve}, with Problems~\ref{pr:vector-d} and~\ref{pr:fibres} and \(T = BC\,r^{nw/2}\).
\end{description}

\subsection*{\S\ref{sec:induction}}
\begin{description}
\item[\ref{pr:gm}] The exact sequence of \(\G_m \subset \A^1\) and the point \(0\).
\item[\ref{pr:dim0}] Frobenius permutes cyclically the geometric points over a closed point of degree \(m\), and its \(m\)-th power acts on each stalk as the local Frobenius.
\item[\ref{pr:extension}] Compare the eigenvalues of \(\Fr\) and of \(\Fr^e\).
\item[\ref{pr:devissage}] The long exact sequence of \(0 \to j_!j^*\cF \to \cF \to i_*i^*\cF \to 0\).
\item[\ref{pr:lisse-open}] Cohomology does not see nilpotents.  Remove the intersections of the irreducible components, then use constructibility and generic smoothness.
\item[\ref{pr:traces}] A point \(u \in U(\F_{q^n})\) lies over a closed point of degree \(m \mid n\), and \(\Fr_u\) is the \((n/m)\)-th power of its local Frobenius.
\item[\ref{pr:higher}] Problem~\ref{so-pr:gysin}.
\item[\ref{pr:purity}] Apply (i) to \(\cF^\vee\), which is pure of weight \(-w\), and use Poincar\'e duality.
\end{description}

\begin{thebibliography}{SGA4\(\frac12\)}

\bibitem[Ald]{Ald}
\emph{A short proof of Deligne's bound for Frobenius eigenvalues, after Alpöge}, contributed by F.~Alderliefste, deposit hexagon:2611.00789v1, 21 November 2026, \url{https://hexagonmath.org/2611.00789v1}.

\bibitem[Alp]{Alp}
L.~Alp\"oge, untitled note announcing a proof of the Riemann hypothesis, 20 November 2026, \url{https://alpo.ge/remark.pdf}.

\bibitem[Del74]{Del74}
P.~Deligne, \emph{La conjecture de Weil.~I}, Publ. Math. IH\'ES \textbf{43} (1974), 273--307.

\bibitem[Del80]{Del80}
P.~Deligne, \emph{La conjecture de Weil.~II}, Publ. Math. IH\'ES \textbf{52} (1980), 137--252.

\bibitem[Mil]{Mil}
J.~S.~Milne, \emph{\'Etale Cohomology}, Princeton, 1980.

\bibitem[SGA4\(\frac12\)]{SGA45}
P.~Deligne, \emph{Cohomologie \'etale} (SGA~4\(\frac12\)), Lecture Notes in Math. \textbf{569}, Springer, 1977.

\bibitem[Zh]{Zh}
D.~Zhang, \emph{On an elementary proof of Weil~II}, arXiv:2611.30749, 23 November 2026 (listed as \emph{Notes on an elementary proof of Weil~II from Hexagon}), \url{https://arxiv.org/abs/2611.30749}.

\end{thebibliography}

\end{document}
